\mxnewpage\section*{EV2-T2 -- Blanc 1}\label{EV2-T2}
\addcontentsline{toc}{section}{EV2-T2 -- Blanc 1}
\textbf{Sujet d’entraînement}. Chapitres 1 à 7.\par Durée : 2 h. Sans calculatrice. Note sur 20. 14–16 janvier 2027.
\subsection*{Partie 1 -- Automatismes (6 points)}
Douze questions, \num{0,5} point chacune. Une seule réponse correcte ; aucune justification demandée. Une réponse incorrecte ou absente rapporte 0, sans retrait de points.
\begin{enumerate}
\item Calculer $3+2\times3$.\par A. $12$\quad B. $8$\quad C. $15$\quad D. $9$
\item Calculer $\frac12+\frac13$.\par A. $\frac{2}{5}$\quad B. $\frac{1}{6}$\quad C. $\frac{5}{3}$\quad D. $\frac{5}{6}$
\item Calculer $\dfrac{2^3}{2^{1}}$.\par A. $4$\quad B. $2$\quad C. $8$\quad D. $1$
\item Une valeur de 50 augmente de \pc{20}. Valeur finale ?\par A. $40$\quad B. $70$\quad C. $60$\quad D. $10$
\item Une valeur de 80 est obtenue après une baisse de \pc{20}. Valeur initiale ?\par A. $64$\quad B. $96$\quad C. $100$\quad D. $60$
\item Hausse de \pc{25} puis baisse de \pc{20} : coefficient global ?\par A. $\frac{9}{10}$\quad B. $\frac{21}{20}$\quad C. $\frac{4}{5}$\quad D. $1$
\item La solution de $2x+3=0$ vaut\par A. $6$\quad B. $-6$\quad C. $\frac{3}{2}$\quad D. $\frac{-3}{2}$
\item La pente entre $\pt{0;3}$ et $\pt{2;9}$ vaut\par A. $-3$\quad B. $2$\quad C. $6$\quad D. $3$
\item La moyenne de $3, 5, 7$ vaut\par A. $3$\quad B. $5$\quad C. $6$\quad D. $15$
\item Un secteur de $90^\circ$ représente quelle fraction du total ?\par A. $\frac{1}{4}$\quad B. $\frac{1}{3}$\quad C. $\frac{3}{4}$\quad D. $\frac{1}{2}$
\item Une probabilité vaut $\dfrac{2}{5}$. Son complément vaut\par A. $\frac{2}{5}$\quad B. $\frac{5}{2}$\quad C. $\frac{1}{5}$\quad D. $\frac{3}{5}$
\item Pour $f(x)=x^2+3$, $f(-2)$ vaut\par A. $12$\quad B. $7$\quad C. $-1$\quad D. $5$
\end{enumerate}
\mxnewpage\subsection*{Partie 2 -- Exercices indépendants (14 points)}
Toute démarche doit être justifiée. Les résultats exacts sont attendus sauf indication contraire.
\mxneed{12}
\subsection*{Exercice 1 -- Deux tirages (5 points)}
Une urne contient quatre boules rouges et deux boules bleues, indiscernables au toucher. On tire successivement deux boules au hasard, sans remettre la première. À chaque tirage, les boules encore présentes sont équiprobables. La troisième question propose un autre protocole, avec remise.
\begin{enumerate}
\item Calculer la probabilité d’obtenir deux boules rouges. \hfill(1 pt)
\item Calculer la probabilité d’obtenir une boule de chaque couleur, en tenant compte des deux ordres possibles. \hfill(2 pts)
\item On remet désormais la première boule avant le second tirage. Calculer la probabilité d’obtenir une boule de chaque couleur et la comparer à celle du protocole sans remise. \hfill(2 pts)
\end{enumerate}
\mxneed{12}
\subsection*{Exercice 2 -- Une boîte (5 points)}
On découpe des carrés de côté $x$ cm dans les quatre coins d’une feuille carrée de 10 cm de côté. En relevant les bords, on forme une boîte sans couvercle. Son volume, en centimètres cubes, est $V(x)=x(10-2x)^2$. On étudie ce modèle sur $\intff{0}{5}$, bornes aplaties comprises.
\begin{enumerate}
\item Calculer $V\prime$ sous forme factorisée. \hfill(1 pt)
\item Étudier les variations. \hfill(2 pts)
\item Déterminer le volume maximal exact et justifier que les bornes ne donnent pas le maximum. \hfill(2 pts)
\end{enumerate}
\mxneed{12}
\subsection*{Exercice 3 -- Reconnaissance et intersection (4 points)}
Dans un repère orthonormé, on considère l’ensemble $\mathcal C$ des points $M\pt{x;y}$ qui vérifient $x^2+y^2+4x-2y-4=0$. Après l’avoir identifié, on étudiera ses intersections avec une droite et une tangente.
\begin{enumerate}
\item Montrer que $\mathcal C$ est un cercle et donner son centre et son rayon. \hfill(1 pt)
\item Déterminer tous les points communs à $\mathcal C$ et à la droite $x=-2$. \hfill(1 pt)
\item Écrire l’équation de la tangente au cercle au point d’intersection de plus grande ordonnée. Justifier sa direction. \hfill(2 pts)
\end{enumerate}
